\begin{table}[H]
\centering
\caption{\textbf{Three measures of reach and their dependence on the adjustment method.}
Each entry is the percentage difference under narrower-scope relative to broader-scope
publication, with its journal-clustered analytic 95\% interval. Weighted estimates use the
saved propensity weights alone; outcome-model estimates standardize the fitted
journal-group-specific models to the common covariate distribution and their intervals
condition on those fitted predictions, so they are descriptive rather than a claim of greater
precision; augmented estimates combine both and are the values reported in the main text. The
final row is a conditional mean and compares the papers that received such citations under each
journal group, not one fixed population. All rows use the separately regenerated comparison
sample of 3,827,491 papers from 20,215 journals.}
\label{tab:decomposition}
\footnotesize
\setlength{\tabcolsep}{5pt}
\begin{tabular}{lrrr}
\toprule
Measure & Weighted & Outcome model & Augmented \\
\midrule
Cited by any other research area & -9.7 (-19.4 to 1.2) & -0.4 & -1.0 (-2.8 to 0.9) \\
Number of other named research areas citing & -16.3 (-26.4 to -4.8) & -5.5 (-6.5 to -4.5) & -5.2 (-8.2 to -2.2) \\
Citations from other research areas & -25.2 (-38.0 to -9.8) & -8.6 & -12.7 (-21.7 to -2.7) \\
Citations from other areas, among papers receiving any & -17.2 (-27.5 to -5.5) & -8.2 & -11.8 (-20.7 to -1.9) \\
\midrule
\multicolumn{4}{l}{\textit{Decomposition of the citation shortfall}} \\
\quad Entry component & \multicolumn{3}{r}{-1.0 (-2.8 to 0.9)} \\
\quad Intensity component & \multicolumn{3}{r}{-11.8 (-20.7 to -1.9)} \\
\quad Intensity minus entry & \multicolumn{3}{r}{-11.0 (-20.0 to -0.9)} \\
\bottomrule
\end{tabular}
\end{table}
